\section{Stinespring dilation}
\label{sec:stinespring}
%% =========================================================

For a normalized Kraus family, the Stinespring construction realizes the
channel through an isometry $V$ into a larger space,
$T(\rho)=\tr_E[V\rho V^\dagger]$; tracing out the environment $E$ returns the
channel.
% A tensor-network diagram for T(\rho)=\tr_E[V\rho V^\dagger] accompanied
% this formula in TNLean (Wolf, Sec. 2.2, proof of Thm. 2.5). QICLean does
% not carry the tenkz diagram pipeline; the equation above is the complete
% statement. See interface_edges.md / the packaging report.

\begin{definition}[Stinespring matrix]\label{def:stinespring_v}
    \lean{stinespringV}
    \leanok
    Given possibly rectangular Kraus operators
    $K_j:\C^{d_{\mathrm{in}}}\to\C^{d_{\mathrm{out}}}$, the
    \emph{Stinespring matrix}
    $V:\C^{d_{\mathrm{in}}}\to
       \C^{d_{\mathrm{out}}}\otimes\C^r$ is defined by
    \begin{align}
        V_{(i,j),k}
        &=(K_j)_{ik}.
        \label{eq:representations_stinespring_entries}
    \end{align}
    Thus $V=\sum_j K_j\otimes\ket{j}$. It is an isometry precisely when
    the Kraus family satisfies the trace-preserving normalization.
\end{definition}

\begin{lemma}[Entry formula for the Stinespring matrix]\label{lem:stinespring_v_apply}
    \lean{stinespringV_apply}
    \leanok
    \uses{def:stinespring_v}
    For output index $i$, input index $k$, and $j\in\{0,\ldots,r-1\}$,
    $V_{(i,j),k}=(K_j)_{ik}$.
\end{lemma}

\begin{definition}[Stinespring matrix with a finite index set]
    \label{def:stinespring_v_gen}
    \lean{stinespringVGen}
    \leanok
    \uses{def:stinespring_v}
    For a Kraus family $(K_j)_{j\in J}$ indexed by an arbitrary finite set
    $J$, define $(V_J)_{(i,j),k}=(K_j)_{ik}$.
\end{definition}

\begin{lemma}[Entry formula for the finite-index Stinespring matrix]
    \label{lem:stinespring_v_gen_apply}
    \lean{stinespringVGen_apply}
    \leanok
    \uses{def:stinespring_v_gen}
    The finite-index Stinespring matrix satisfies
    $(V_J)_{(i,j),k}=(K_j)_{ik}$.
\end{lemma}

\begin{theorem}[Stinespring Gram identity]\label{thm:stinespring_v_conjTranspose_mul}
    \lean{stinespringV_conjTranspose_mul}
    \leanok
    \uses{def:stinespring_v}
    The Stinespring matrix satisfies
    \begin{align}
        V^\dagger V
        &=\sum_{j=0}^{r-1}K_j^\dagger K_j.
        \label{eq:representations_stinespring_gram}
    \end{align}
\end{theorem}

\begin{definition}[Local Stinespring matrix]\label{def:local_stinespring_v}
    \lean{localStinespringV}
    \leanok
    \uses{def:stinespring_v}
    If a Stinespring matrix $V:\C^{d_{\mathrm{in}}}\to
    \C^{d_{\mathrm{out}}}$ acts on the first factor while a finite-dimensional
    space $R$ is left unchanged, the corresponding local Stinespring matrix is
    $W=V\otimes\Id_R$.
\end{definition}

\begin{theorem}[Local Stinespring isometry]
    \label{thm:local_stinespring_v_conjTranspose_mul}
    \lean{localStinespringV_conjTranspose_mul}
    \leanok
    \uses{def:local_stinespring_v}
    If $V^\dagger V=\Id$, then
    $(V\otimes\Id_R)^\dagger(V\otimes\Id_R)=\Id$.
\end{theorem}

\begin{proof}\leanok
    Use $(A\otimes B)(C\otimes D)=AC\otimes BD$ and
    $\Id\otimes\Id=\Id$.
\end{proof}

\begin{theorem}[Stinespring dual representation]\label{thm:stinespring_dual_representation}
    \lean{stinespring_dual_representation}
    \leanok
    \uses{def:stinespring_v, def:trace_pairing_adjoint}
    Let $K_j:\C^{d}\to\C^{d'}$, $j=0,\dots,r-1$, be an arbitrary family of
    matrices and let $V=\sum_j K_j\otimes\ket{j}$. Then, for every observable
    $A\in\MN{d'}$ on the output space,
    \begin{align}
        V^\dagger(A\otimes\Id_r)V
        &=\sum_j K_j^\dagger A K_j,
        \label{eq:representations_stinespring_dual}
    \end{align}
    an identity in $\MN{d}$ on the input space, valid for arbitrary $K$
    independently of any channel. When the $K_j$ are Kraus operators of a map
    $T:\MN{d}\to\MN{d'}$, the right-hand side is the dual
    (Heisenberg-picture) map $T^*(A)$ of \ref{def:trace_pairing_adjoint}; the
    square case is $d=d'$.
\end{theorem}

\begin{proof}
    \leanok
    \uses{def:stinespring_v}
    Compute $(V^\dagger(A\otimes\Id_r)V)_{ab}$
    by summing over the product index $(i,j)$;
    the $\delta_{jl}$ from $\Id_r$ collapses the sum over $l$, giving
    $\sum_j K_j^\dagger A K_j$. For Kraus operators of $T$, the trace
    pairing identifies this sum with $T^*(A)$, as displayed in the proof of
    Theorem~\ref{thm:stinespring_rectangular}.
\end{proof}

\begin{theorem}[The Stinespring isometry condition iff trace preservation]\label{thm:stinespringV_isometry_iff_kraus_normalized}
    \lean{stinespringV_isometry_iff_kraus_normalized}
    \leanok
    \uses{def:stinespring_v, def:trace_preserving}
    In the square specialization
    $d_{\mathrm{in}}=d_{\mathrm{out}}=D$, one has
    $V^\dagger V=\Id_D$ if and only if
    $\sum_j K_j^\dagger K_j=\Id_D$.
    In particular, $V$ is an isometry precisely when the Kraus map is
    trace-preserving.
\end{theorem}

\begin{proof}
    \leanok
    \uses{def:stinespring_v}
    One has
    $(V^\dagger V)_{ab}
    =\sum_{(i,j)}\overline{V_{(i,j),a}}V_{(i,j),b}
    =\sum_j\sum_i\overline{(K_j)_{ia}}(K_j)_{ib}
    =(\sum_jK_j^\dagger K_j)_{ab}$.
\end{proof}

\begin{theorem}[Stinespring Schr\"odinger representation]\label{thm:stinespring_schrodinger_representation}
    \lean{stinespring_schrodinger_representation}
    \leanok
    \uses{def:stinespring_v, def:cp_map, def:partial_trace}
    The Kraus map $T(\rho)=\sum_jK_j\rho K_j^\dagger$ equals
    the partial trace over the dilation space:
    \begin{align}
        T(\rho)_{ij}
        &=\sum_{k=0}^{r-1}(V\rho V^\dagger)_{(i,k),(j,k)}
          =\tr_r(V\rho V^\dagger).
        \label{eq:representations_stinespring_schrodinger}
    \end{align}
    This is the Schr\"odinger-picture Stinespring representation.
\end{theorem}

\begin{proof}
    \leanok
    \uses{def:stinespring_v}
    Expand the middle expression in
    (\ref{eq:representations_stinespring_schrodinger}) and sum over $k$ to
    recover $\sum_kK_k\rho K_k^\dagger$.
\end{proof}

\begin{theorem}[Existence of a Stinespring dilation]
    \label{thm:exists_stinespring_dilation}
    \lean{exists_stinespring_dilation}
    \leanok
    \uses{def:cp_map, def:stinespring_v}
    Every completely positive map admits an ancilla dimension $r$, Kraus operators
    $\{K_j\}_{j=0}^{r-1}$, and the concrete representation
    $\pi(A)=A\otimes\Id_r$ such that $E(A)=V^\dagger\pi(A)V$.
\end{theorem}

\begin{proof}\leanok
    Choose a Kraus representation of $E$ and apply the explicit Heisenberg-picture
    Stinespring formula.
\end{proof}

\begin{theorem}[CPTP maps admit isometric Stinespring dilations]
    \label{thm:exists_stinespring_isometry_of_cptp}
    \lean{exists_stinespring_isometry_of_cptp}
    \leanok
    \uses{def:channel, def:stinespring_v}
    Every quantum channel admits a Stinespring dilation whose matrix $V$ is an
    isometry, equivalently $V^\dagger V=\Id$.
\end{theorem}

\begin{proof}\leanok
    Choose Kraus operators for the channel and use trace preservation to obtain
    the Stinespring isometry condition $V^\dagger V=\Id$.
\end{proof}

\begin{theorem}[Stinespring representation, rectangular form]
    \label{thm:stinespring_rectangular}
    \lean{ChoiRectangular.exists_stinespringV_of_isKrausCP,
        ChoiRectangular.exists_stinespringV_schrodinger_of_isKrausCP,
        ChoiRectangular.exists_stinespringV_schrodinger_of_isKrausCPTP,
        ChoiRectangular.exists_stinespringV_choiRank_of_isKrausCP,
        ChoiRectangular.exists_stinespringV_pairing_of_isKrausCP,
        ChoiRectangular.exists_stinespringV_pairing_choiRank_of_isKrausCP}
    \leanok
    \uses{def:is_kraus_cp, def:trace_pairing_adjoint, def:kraus_rank_rect}
    Let $T:\MN{d}\to\MN{d'}$ be completely positive with Choi matrix $\tau$ and
    dual $T^*$. Then for every $r\geq\rank(\tau)$ there is a
    $V:\C^{d}\to\C^{d'}\otimes\C^{r}$ such that, for all $A\in\MN{d'}$,
    \begin{align}
        T^*(A)&=V^\dagger(A\otimes\Id_r)V,
        \label{eq:stinespring_rectangular}
    \end{align}
    and $V$ is an isometry, $V^\dagger V=\Id_d$, if and only if $T$ is trace
    preserving. In the Schr\"odinger picture one may choose the same type of
    dilation so that
    \begin{align}
        T(\rho)&=\tr_{\C^r}(V\rho V^\dagger)
        \label{eq:stinespring_rectangular_schrodinger}
    \end{align}
    for every $\rho\in\MN{d}$. In particular $r=\rank(\tau)$ is an admissible
    ancilla dimension; by
    Lemma~\ref{lem:stinespring_rectangular_least_dim} no smaller
    ancilla dimension is admissible, and a dilation with $r=\rank(\tau)$ is
    called minimal. This is \cite[Theorem~2.2]{Wolf2012Quantum}.
\end{theorem}

\begin{proof}
    \leanok
    \uses{def:kraus_rank_rect, def:stinespring_v, thm:kraus_rect_rank_minimal,
        thm:kraus_rect_normalization, thm:stinespring_dual_representation,
        thm:stinespring_v_conjTranspose_mul,
        thm:trace_pairing_adjoint_identity, thm:trace_nondeg}
    By minimality of the Kraus rank, $T$ has a Kraus family of exactly
    $\rank(\tau)$ operators; padding it with zeros gives a family
    $\{K_j\}_{j=0}^{r-1}$ for every $r\geq\rank(\tau)$. Set
    $V:=\sum_j K_j\otimes\ket{j}$ for an orthonormal basis $\{\ket{j}\}$ of
    $\C^{r}$, so that $V^\dagger(A\otimes\Id_r)V=\sum_j K_j^\dagger AK_j$. The
    trace pairing gives
    \begin{align}
        \tr\bigl[T^*(A)X\bigr]
        &=\tr\bigl[A\,T(X)\bigr]
          =\tr\Bigl[A\sum_j K_jXK_j^\dagger\Bigr]
          =\tr\Bigl[\sum_j K_j^\dagger AK_j\,X\Bigr]
        \label{eq:stinespring_rectangular_pairing}
    \end{align}
    for every $X\in\MN{d}$, and nondegeneracy of the trace pairing yields
    (\ref{eq:stinespring_rectangular}). Finally
    $V^\dagger V=\sum_j K_j^\dagger K_j$, and the latter equals $\Id_d$
    exactly when $T$ is trace preserving.
\end{proof}

\begin{lemma}[Least dilation dimension]\label{lem:stinespring_rectangular_least_dim}
    \lean{ChoiRectangular.choiRank_le_of_stinespring_dual_representation}
    \leanok
    \uses{def:trace_pairing_adjoint, def:kraus_rank_rect}
    Let $T:\MN{d}\to\MN{d'}$ be a linear map with Choi matrix $\tau$ and dual
    $T^*$. If for some ancilla dimension $r$ there is a
    $V:\C^{d}\to\C^{d'}\otimes\C^{r}$ such that
    (\ref{eq:stinespring_rectangular}) holds for all $A\in\MN{d'}$, then
    $\rank(\tau)\le r$. Together with
    Theorem~\ref{thm:stinespring_rectangular}, the least admissible ancilla
    dimension is $\rank(\tau)$, and dilations with $r=\rank(\tau)$ are
    minimal. This is the discussion following
    \cite[Theorem~2.2]{Wolf2012Quantum}.
\end{lemma}

\begin{proof}
    \leanok
    \uses{def:stinespring_v, thm:kraus_rect_rank_minimal,
        thm:stinespring_dual_representation,
        thm:trace_pairing_adjoint_identity, thm:trace_nondeg}
    The Kraus operators are recovered from $V$ as the blocks
    $K_j=(\Id_{d'}\otimes\bra{j})V$, which satisfy
    $V=\sum_jK_j\otimes\ket{j}$, so that
    $T^*(A)=\sum_jK_j^\dagger AK_j$ for every $A\in\MN{d'}$. For every
    $X\in\MN{d}$, the trace pairing gives
    \begin{align}
        \tr\bigl[A\,T(X)\bigr]
        &=\tr\bigl[T^*(A)X\bigr]
          =\tr\Bigl[\sum_jK_j^\dagger AK_j\,X\Bigr]
          =\tr\Bigl[A\sum_jK_jXK_j^\dagger\Bigr],
        \label{eq:stinespring_rectangular_least_pairing}
    \end{align}
    and nondegeneracy of the trace pairing yields
    $T(X)=\sum_jK_jXK_j^\dagger$: an $r$-operator Kraus family of $T$. Any
    $r$-operator Kraus family forces $\rank(\tau)\le r$, the bound in the
    minimality of the Kraus rank.
\end{proof}

%% =========================================================
\section{Ordered CP maps, Radon--Nikodym, and open-system representation}
\label{sec:ordered_cp}
%% =========================================================

The next three results form the structural part of~\cite[Theorems~2.3--2.5]{Wolf2012Quantum}:
the ordered CP-map theorem, the Radon--Nikodym theorem for completely positive
maps, and the open-system representation. They refine the Stinespring dilation
by tracking how two CP maps related by domination or decomposition sit inside
a common dilation space.

\begin{definition}[CP partial order]\label{def:cp_dominates}
    \lean{CPDominates, CPDominates.refl}
    \leanok
    \uses{def:cp_map}
    For linear maps $S,T:\MN{d_{\rm in}}\to\MN{d_{\rm out}}$, write
    $T\le S$ (in the CP order) when $S-T$ is completely positive.
\end{definition}

\begin{definition}[Wolf auxiliary matrix]\label{def:ordered_cp_stinespring_w}
    \lean{stinespringW, stinespringW_apply}
    \leanok
    \uses{def:stinespring_v}
    Let $V:\C^d\to\C^{d'}\otimes\C^r$ be a supplied Stinespring
    matrix, and let
    $\ket{\Omega}=d'^{-1/2}\sum_{a=1}^{d'}\ket{a,a}$ be normalized.
    Wolf's auxiliary operator is
    \begin{align}
        W
        &:=(\Id_r\otimes\bra{\Omega})(V\otimes\Id_{d'})
          :\C^d\otimes\C^{d'}\longrightarrow\C^r,
        \label{eq:ordered_cp_w_definition}
    \end{align}
    or, in coordinates, $W_{j,(i,a)}=d'^{-1/2}V_{(a,j),i}$.
\end{definition}

\begin{lemma}[The $V$--$W$ correspondence]
    \label{lem:ordered_cp_v_w_correspondence}
    \lean{stinespringW_injective, stinespringW_kronecker_mul,
        stinespring_eq_kronecker_mul_iff_W_eq_mul}
    \leanok
    \uses{def:ordered_cp_stinespring_w}
    If $d'>0$, the assignment $V\mapsto W$ is one-to-one. Moreover, for
    every $C:\C^{r_2}\to\C^{r_1}$,
    \[
        V_1=(\Id_{d'}\otimes C)V_2
        \quad\Longleftrightarrow\quad
        W_1=CW_2.
    \]
\end{lemma}

\begin{proof}\leanok
    The coordinate formula
    $W_{j,(i,a)}=d'^{-1/2}V_{(a,j),i}$ is invertible because $d'>0$.
    Expanding the Kronecker product shows directly that multiplication by
    $\Id_{d'}\otimes C$ on $V$ becomes multiplication by $C$ on $W$.
\end{proof}

\begin{lemma}[The auxiliary Gram matrix is the Choi matrix]
    \label{lem:ordered_cp_stinespring_w_gram}
    \lean{stinespringW_conjTranspose_mul_self}
    \leanok
    \uses{def:ordered_cp_stinespring_w, def:choi_matrix_rect,
        thm:stinespring_dual_representation}
    Suppose that $T:\MN{d'}\to\MN{d}$ has the supplied representation
    $T(A)=V^\dagger(A\otimes\Id_r)V$. If $\tau$ is the normalized Choi
    matrix of $T$, then
    \begin{align}
        W^\dagger W=\tau.
        \label{eq:ordered_cp_w_gram}
    \end{align}
\end{lemma}

\begin{proof}\leanok
    Write $V=\sum_jK_j\otimes\ket j$. Then
    $T(A)=\sum_jK_j^\dagger A K_j$, so $T$ has Kraus operators
    $K_j^\dagger$. Both sides of
    (\ref{eq:ordered_cp_w_gram}) have entries
    \[
        \frac1{d'}\sum_j
        \overline{(K_j)_{a i}}(K_j)_{b k}
    \]
    at the pair of indices $(i,a),(k,b)$.
\end{proof}

\begin{lemma}[Rectangular contraction factorization]
    \label{lem:ordered_cp_rectangular_contraction_factorization}
    \lean{Matrix.exists_contraction_mul_of_conjTranspose_mul_le,
        Matrix.sqNorm_mulVec_le_of_conjTranspose_mul_le,
        Matrix.exists_contraction_mul_of_sqNorm_le}
    \leanok
    Let $W_i:H\to R_i$ be finite-dimensional complex matrices. If
    \begin{align}
        \lVert W_1\psi\rVert^2
        \leq \lVert W_2\psi\rVert^2
        \qquad(\psi\in H),
        \label{eq:ordered_cp_norm_comparison}
    \end{align}
    then there is a contraction $C:R_2\to R_1$ such that
    $W_1=CW_2$.
\end{lemma}

\begin{proof}\leanok
    Equivalently, (\ref{eq:ordered_cp_norm_comparison}) is
    $W_1^\dagger W_1\leq W_2^\dagger W_2$. Take the positive square root
    $S$ of $W_2^\dagger W_2-W_1^\dagger W_1$. After adjoining zero rows,
    the matrices with row blocks $(W_1,S,0)$ and $(0,0,W_2)$ have the
    same Gram matrix. A unitary carrying the second matrix to the first
    has an $R_1\times R_2$ corner $C$ satisfying $W_1=CW_2$; orthonormality
    of its columns gives $C^\dagger C\leq\Id_{R_2}$.
\end{proof}

\begin{lemma}[Minimality makes the auxiliary matrix surjective]
    \label{lem:ordered_cp_minimal_w_surjective}
    \lean{stinespringW_mulVec_surjective_of_minimal,
        Matrix.eq_of_mul_eq_mul_of_mulVec_surjective}
    \leanok
    \uses{lem:ordered_cp_stinespring_w_gram,
        lem:stinespring_rectangular_least_dim}
    If the supplied dominating dilation is minimal,
    $r_2=\rank(\tau_2)$, then $W_2$ is surjective. Consequently a
    factorization $W_1=CW_2$ determines $C$ uniquely.
\end{lemma}

\begin{proof}\leanok
    The Gram identity gives
    $\rank(W_2)=\rank(W_2^\dagger W_2)=\rank(\tau_2)=r_2$, so $W_2$ has
    full row rank and is surjective. A surjective right factor can be
    cancelled.
\end{proof}

\begin{theorem}[Relation between ordered CP maps]
    \label{thm:cp_dominates_stinespring_contraction}
    \lean{CPDominates.exists_supplied_stinespring_contraction}
    \leanok
    \uses{def:cp_dominates, def:ordered_cp_stinespring_w,
        lem:ordered_cp_stinespring_w_gram,
        lem:ordered_cp_rectangular_contraction_factorization,
        lem:ordered_cp_v_w_correspondence,
        lem:ordered_cp_minimal_w_surjective}
    Let $T_i:\MN{d'}\to\MN{d}$ be completely positive maps with
    $T_1\leq T_2$. Suppose that, for possibly distinct ancilla dimensions
    $r_i$, matrices $V_i:\C^d\to\C^{d'}\otimes\C^{r_i}$ are supplied and
    satisfy
    \[
        T_i(A)=V_i^\dagger(A\otimes\Id_{r_i})V_i
        \qquad(A\in\MN{d'}).
    \]
    Then there is a contraction $C:\C^{r_2}\to\C^{r_1}$ such that
    \begin{align}
        V_1=(\Id_{d'}\otimes C)V_2.
        \label{eq:representations_ordered_intertwining}
    \end{align}
    If $V_2$ is minimal, in the source sense
    $r_2=\rank(\tau_2)$, then $C$ is unique. This is
    \cite[Theorem~2.3 and Equation~(2.13)]{Wolf2012Quantum}.
\end{theorem}

\begin{proof}\leanok
    Complete positivity of $T_2-T_1$ and the Choi correspondence give
    $\tau_1\leq\tau_2$. By
    Lemma~\ref{lem:ordered_cp_stinespring_w_gram},
    \[
        W_1^\dagger W_1=\tau_1\leq\tau_2=W_2^\dagger W_2,
    \]
    which is precisely the squared norm comparison
    (\ref{eq:ordered_cp_norm_comparison}). The rectangular factorization
    lemma gives a contraction $C$ with $W_1=CW_2$. Inverting
    (\ref{eq:ordered_cp_w_definition}) gives
    (\ref{eq:representations_ordered_intertwining}) with the same supplied
    $V_i$.

    If $r_2=\rank(\tau_2)$, then
    $\rank(W_2)=\rank(W_2^\dagger W_2)=\rank(\tau_2)=r_2$; hence $W_2$ is
    surjective. Thus two coefficients satisfying $W_1=CW_2=C'W_2$ agree,
    proving uniqueness.
\end{proof}

\begin{definition}[Block-top contraction]\label{def:block_top_rows}
    \lean{Matrix.blockTopRows}
    \leanok
    For natural numbers $r,s$, let $C=C_{r,s}\in\C^{r\times(r+s)}$ be the
    rectangular $r\times(r+s)$ matrix whose rows are the first $r$ rows
    of the identity on $\C^{r+s}$: $C_{ij}=1$ if $j=i<r$ and $0$
    otherwise.
\end{definition}

\begin{lemma}[$C$ is a co-isometry]\label{lem:block_top_rows_co_isometry}
    \lean{Matrix.blockTopRows_mul_conjTranspose}
    \leanok
    \uses{def:block_top_rows}
    The block-top matrix satisfies $CC^\dagger=\Id_r$.
\end{lemma}

\begin{proof}\leanok
    Direct entrywise computation:
    $(CC^\dagger)_{ii'}=\sum_jC_{ij}\overline{C_{i'j}}$
    collapses to $\delta_{ii'}$.
\end{proof}

\begin{lemma}[$C$ is a contraction]\label{lem:block_top_rows_contraction}
    \lean{Matrix.blockTopRows_conjTranspose_mul_le_one}
    \leanok
    \uses{def:block_top_rows}
    The block-top matrix satisfies $C^\dagger C\le\Id_{r+s}$.
\end{lemma}

\begin{proof}\leanok
    The matrix $C^\dagger C$ is diagonal with a $1$ on the first $r$
    coordinates and a $0$ on the last $s$, so $\Id-C^\dagger C$ is
    positive semidefinite.
\end{proof}

\begin{lemma}[Intertwining of Stinespring isometries]\label{lem:stinespring_intertwine_append}
    \lean{stinespringV_eq_kronecker_blockTopRows_mul_append}
    \leanok
    \uses{def:stinespring_v, def:block_top_rows}
    Let $K:\{0,\ldots,r-1\}\to\MN{D}$ and
    $L:\{0,\ldots,s-1\}\to\MN{D}$ be two Kraus families. Write
    $K\mathbin{+\!+}L$ for their concatenation as a family indexed by
    $\{0,\ldots,r+s-1\}$. Then
    \begin{align}
        V_K
        &=(\Id_D\otimes C_{r,s})V_{K+\!+L}.
        \label{eq:representations_stinespring_intertwining}
    \end{align}
\end{lemma}

\begin{proof}\leanok
    Entrywise expansion of both sides of
    (\ref{eq:representations_stinespring_intertwining}): the Kronecker product
    with $C_{r,s}$ picks out the first $r$ coordinates of the dilation space,
    matching the action of $V_{K+\!+L}$ on those coordinates, which is
    precisely $V_K$.
\end{proof}

\begin{corollary}[Canonical square realization]\label{cor:cp_dominates_stinespring_contraction_canonical}
    \lean{CPDominates.exists_stinespring_contraction}
    \leanok
    \uses{def:cp_dominates, def:stinespring_v, def:block_top_rows,
        lem:block_top_rows_contraction, lem:stinespring_intertwine_append}
    Let $T_1,T_2:\MN{D}\to\MN{D}$ be CP maps with $T_1\le T_2$.
    Then there exist ancilla dimensions $r_1,m$, Heisenberg-form Stinespring
    matrices $V_1:\C^D\to\C^D\otimes\C^{r_1}$ and
    $V_2:\C^D\to\C^D\otimes\C^m$ realizing $T_1,T_2$ via
    $T_i(A)=V_i^\dagger(A\otimes\Id)V_i$, and a rectangular contraction
    $\widetilde C:\C^m\to\C^{r_1}$ with
    $\widetilde C^\dagger\widetilde C\le\Id_m$ such that
    \begin{align}
        V_1
        &=(\Id_D\otimes\widetilde C)V_2.
        \label{eq:representations_ordered_intertwining_canonical}
    \end{align}
    This is the explicit square corollary obtained by constructing both
    dilation matrices, rather than the supplied-dilation statement of
    Theorem~\ref{thm:cp_dominates_stinespring_contraction}.
\end{corollary}

\begin{proof}\leanok
    \uses{thm:exists_stinespring_dilation, lem:block_top_rows_contraction,
        lem:stinespring_intertwine_append}
    Choose a Heisenberg-form Kraus family $K_1$ of $T_1$ with Stinespring
    matrix $V_1=V_{K_1}$, and Kraus operators $L$ of the CP map
    $T_2-T_1$ in Schr\"odinger orientation. Let $L^\dagger$ denote the
    conjugate-transposed family. Set $m=r_1+s$ and let
    $K_2=K_1\mathbin{+\!+}L^\dagger$. Then
    \begin{align}
        T_2(A)
        &=T_1(A)+(T_2-T_1)(A) \notag\\
        &=\sum_i(K_1^i)^\dagger A K_1^i+\sum_jL_j A L_j^\dagger \notag\\
        &=\sum_i(K_1^i)^\dagger A K_1^i
          +\sum_j(L_j^\dagger)^\dagger A L_j^\dagger.
        \notag
    \end{align}
    This is the Heisenberg form for $K_2$, yielding
    $T_2(A)=V_{K_2}^\dagger(A\otimes\Id)V_{K_2}$. Taking
    $\widetilde C=C_{r_1,s}$, the intertwining
    $V_{K_1}=(\Id_D\otimes\widetilde C)V_{K_2}$ is
    Lemma~\ref{lem:stinespring_intertwine_append}, and
    $\widetilde C^\dagger\widetilde C\le\Id$ is
    Lemma~\ref{lem:block_top_rows_contraction}.
\end{proof}

\begin{definition}[Block-diagonal projectors on the dilation space]
    \label{def:block_diag_projectors}
    \lean{Matrix.blockDiagTopProj, Matrix.blockDiagBotProj}
    \leanok
    \uses{def:block_top_rows, lem:block_top_rows_contraction}
    For $r,s$, let $P_{\mathrm{top}}=C_{r,s}^\dagger C_{r,s}$ and
    $P_{\mathrm{bot}}=\Id-P_{\mathrm{top}}$. Both are PSD
    (for $P_{\mathrm{bot}}$, by Lemma~\ref{lem:block_top_rows_contraction})
    and $P_{\mathrm{top}}+P_{\mathrm{bot}}=\Id$.
\end{definition}

\begin{theorem}[Radon--Nikodym theorem for finite quantum instruments]
    \label{thm:cp_radon_nikodym_of_stinespring}
    \lean{IsCPMap.radon_nikodym_of_stinespring}
    \leanok
    \uses{def:cp_map, def:stinespring_v}
    Let $\{T_i\}_{i\in I}$ be a nonempty finite family of completely positive
    maps, let $T:\MN{D}\to\MN{D}$ satisfy $\sum_{i\in I}T_i=T$, and
    suppose that $T$ has a supplied Stinespring representation
    \begin{align}
        T(A)
        &=V^\dagger(A\otimes\Id_r)V.
        \label{eq:representations_radon_dilation}
    \end{align}
    Then there are positive semidefinite operators $P_i\in\MN{r}$ with
    \begin{align}
        \sum_{i\in I}P_i
        &=\Id_r.
        \label{eq:representations_radon_resolution}
    \end{align}
    For every $i\in I$ and $A\in\MN{D}$,
    \begin{align}
        T_i(A)
        &=V^\dagger(A\otimes P_i)V.
        \label{eq:representations_radon_components}
    \end{align}
    This is \cite[Theorem~2.4]{Wolf2012Quantum}, with the nonempty-family condition
    made explicit; see
    \cite{gap:wolf_radon_nikodym_nonempty_family}.
\end{theorem}

\begin{proof}\leanok
    \uses{thm:kraus_rectangular_freedom'}
    Choose Kraus families for the component maps $T_i$ and group all their
    Kraus operators into one finite family. Add zero Kraus operators to one
    component so that this grouped family has at least as many rows as the
    supplied Stinespring dilation. Write the grouped Kraus operators as
    $B_\alpha$, with a label map $\ell(\alpha)\in I$, and write the Kraus
    operators obtained from the blocks of $V$ as $A_j$. Rectangular Kraus
    freedom gives a matrix $W$ with
    \begin{align}
        W^\dagger W
        &=\Id_r, \notag\\
        B_\alpha
        &=\sum_{j=0}^{r-1}W_{\alpha j}A_j.
        \notag
    \end{align}
    For the fibre over $i$, set
    \begin{align}
        (C_i)_{\alpha j}
        &=
        \begin{cases}
            \overline{W_{\alpha j}}, & \ell(\alpha)=i,\\
            0, & \ell(\alpha)\ne i,
        \end{cases}
        \notag\\
        P_i
        &=C_i^\dagger C_i.
        \label{eq:representations_radon_effects}
    \end{align}
    Then $P_i\ge0$, and (\ref{eq:representations_radon_resolution}) follows
    from $W^\dagger W=\Id_r$. Finally,
    \begin{align}
        V^\dagger(A\otimes P_i)V
        &=\sum_{\ell(\alpha)=i}B_\alpha A B_\alpha^\dagger
          =T_i(A),
        \notag
    \end{align}
    which proves the claim.
\end{proof}

\begin{theorem}[Binary Radon--Nikodym theorem for CP maps]
    \label{thm:cp_exists_radon_nikodym}
    \lean{IsCPMap.exists_radon_nikodym}
    \leanok
    \uses{def:cp_map, def:stinespring_v, def:block_diag_projectors}
    Let $T_1,T_2:\MN{D}\to\MN{D}$ be completely positive. There exist
    an ancilla dimension $m$, a Kraus family $K$ with Stinespring matrix
    $V=V_K$, and positive semidefinite operators
    $P_1,P_2:\C^m\to\C^m$ with $P_1+P_2=\Id_m$ such that, for
    $i=1,2$ and every $A\in\MN{D}$,
    \begin{align}
        T_i(A)
        &=V^\dagger(A\otimes P_i)V.
        \label{eq:representations_binary_radon}
    \end{align}
\end{theorem}

\begin{remark}[Relation with Wolf's Radon--Nikodym theorem]
    \label{rem:cp_radon_nikodym_scope}
    % Source-faithful scope note: Wolf's statement is the finite-family
    % theorem relative to a supplied Stinespring dilation.
    Theorem~\ref{thm:cp_radon_nikodym_of_stinespring} treats a finite family
    relative to a supplied Stinespring dilation. The binary case in
    Theorem~\ref{thm:cp_exists_radon_nikodym} instead constructs a common
    dilation from the Kraus families of $T_1$ and $T_2$.
\end{remark}

\begin{proof}\leanok
    \uses{thm:cp_dominates_stinespring_contraction, thm:exists_stinespring_dilation,
        def:block_diag_projectors}
    Take Heisenberg-form Kraus families $K_1$ for $T_1$ and
    Schr\"odinger-form $L$ for $T_2$, and form
    $K=K_1\mathbin{+\!+}L^\dagger$ on $\C^{r_1+s}$. Choose
    $P_1=P_{\mathrm{top}}$, $P_2=P_{\mathrm{bot}}$ as in
    Definition~\ref{def:block_diag_projectors}. The Kronecker identity
    $A\otimes(C^\dagger C)
    =(\Id\otimes C)^\dagger(A\otimes\Id)(\Id\otimes C)$ rewrites
    $V^\dagger(A\otimes P_{\mathrm{top}})V$ as
    $((\Id\otimes C)V)^\dagger(A\otimes\Id)((\Id\otimes C)V)$,
    which equals $V_{K_1}^\dagger(A\otimes\Id)V_{K_1}=T_1(A)$ by
    Lemma~\ref{lem:stinespring_intertwine_append}. For the complementary
    block, $A\otimes P_{\mathrm{bot}}
    =A\otimes\Id-A\otimes P_{\mathrm{top}}$, and sandwiching by $V$ yields
    $(T_1+T_2)(A)-T_1(A)=T_2(A)$.
\end{proof}

\begin{theorem}[Open-system representation of quantum channels]
    \label{thm:channel_exists_stinespring_open_system}
    \lean{IsChannel.exists_stinespring_open_system}
    \leanok
    \uses{def:channel, def:stinespring_v, def:partial_trace}
    Every quantum channel $T:\MN{D}\to\MN{D}$ admits a Stinespring
    isometry $V$ on an ancilla space $\C^r$ such that
    \begin{align}
        T(\rho)_{ij}
        &=\sum_{k=0}^{r-1}(V\rho V^\dagger)_{(i,k),(j,k)}
          =\tr_r(V\rho V^\dagger).
        \label{eq:representations_open_system}
    \end{align}
\end{theorem}

\begin{proof}\leanok
    \uses{thm:exists_stinespring_isometry_of_cptp,
        thm:stinespring_schrodinger_representation}
    Apply the existence of an isometric Stinespring dilation
    (Theorem~\ref{thm:exists_stinespring_isometry_of_cptp}) and the
    Schr\"odinger-picture identity
    (\ref{eq:representations_stinespring_schrodinger}).
\end{proof}

\begin{theorem}[Open-system representation, partial-trace form]
    \label{thm:channel_exists_stinespring_open_system_traceRight}
    \lean{IsChannel.exists_stinespring_open_system_traceRight}
    \leanok
    \uses{def:channel, def:stinespring_v, def:partial_trace}
    Every quantum channel $T$ admits an isometric Stinespring dilation
    $V$ such that $T(\rho)=\tr_E(V\rho V^\dagger)$.
\end{theorem}

\begin{proof}\leanok
    \uses{thm:channel_exists_stinespring_open_system}
    Rewrite the componentwise identity
    (\ref{eq:representations_open_system}) as a partial trace over the
    ancilla factor.
\end{proof}

\begin{definition}[First environment embedding]
    \label{def:first_env_embedding}
    \lean{Matrix.firstEnvEmbedding}
    \leanok
    For $r\geq1$, let $W_0:\C^D\to\C^D\otimes\C^r$ be the isometry
    $W_0x=x\otimes e_0$.
\end{definition}

\begin{lemma}[First environment embedding is an isometry]
    \label{lem:first_env_embedding_conjTranspose_mul_self}
    \lean{Matrix.firstEnvEmbedding_conjTranspose_mul_self}
    \leanok
    \uses{def:first_env_embedding}
    The first-environment embedding satisfies $W_0^\dagger W_0=\Id$.
\end{lemma}

\begin{proof}\leanok
\end{proof}

\begin{theorem}[Open-system representation, unitary form]
    \label{thm:channel_exists_stinespring_open_system_unitary}
    \lean{IsChannel.exists_stinespring_open_system_unitary}
    \leanok
    \uses{def:channel, def:partial_trace, def:first_env_embedding}
    For $D\geq1$, every quantum channel $T$ on $\C^D$ admits an
    environment dimension $r\geq1$ and a unitary
    $U$ on $\C^D\otimes\C^r$ such that, for every matrix
    $\rho$ on $\C^D$,
    \begin{align}
        T(\rho)
        &=\tr_E\!\left(UW_0\rho W_0^\dagger U^\dagger\right).
        \label{eq:representations_open_unitary}
    \end{align}
\end{theorem}

\begin{proof}\leanok
    \uses{thm:channel_exists_stinespring_open_system_traceRight,
        lem:exists_unitary_mul_eq_of_gram_eq,
        lem:first_env_embedding_conjTranspose_mul_self}
    Take the isometric open-system representation
    $T(\rho)=\tr_E(V\rho V^\dagger)$. Since
    $V^\dagger V=\Id=W_0^\dagger W_0$,
    Lemma~\ref{lem:exists_unitary_mul_eq_of_gram_eq}
    gives a unitary $U$ with $V=UW_0$. Substituting this identity into the
    partial trace formula gives (\ref{eq:representations_open_unitary}).
\end{proof}

\begin{theorem}[Open-system representation]
    \label{thm:channel_exists_open_system_unitary_rectangular}
    \lean{Channel.openSystemInputEmbedding,
        Channel.openSystemInput,
        Channel.openSystemInput_eq_embedding,
        Channel.openSystemPartialTrace,
        Channel.output_dimension_pos_of_isKrausCPTP,
        Channel.IsKrausCPTP.exists_openSystem_unitary}
    \leanok
    \uses{def:is_kraus_cptp, def:partial_trace, thm:stinespring_rectangular,
        thm:kraus_rect_rank_minimal, lem:exists_unitary_mul_eq_of_gram_eq}
    Let $d\geq1$ and let $T:\MN{d}\to\MN{d'}$ be completely positive and
    trace-preserving. Then $d'\geq1$, and there are a normalized vector
    $\varphi\in\C^{d'}\otimes\C^{d'}$ and a unitary $U$ on
    $\C^d\otimes\C^{d'}\otimes\C^{d'}$, a space of total dimension
    $d(d')^2$, such that, for every $\rho\in\MN{d}$,
    \begin{align}
        T(\rho)
        &=\tr_{\C^d\otimes\C^{d'}}\!\left[
          U\bigl(\rho\otimes\ketbra{\varphi}{\varphi}\bigr)U^\dagger
          \right].
        \label{eq:representations_open_system_rectangular_unitary}
    \end{align}
    Here the partial trace is over the first two tensor factors
    $\C^d\otimes\C^{d'}$ and retains the final factor $\C^{d'}$. This is
    \cite[Theorem~2.5, Equation~(2.14)]{Wolf2012Quantum}.
\end{theorem}

\begin{proof}\leanok
    \uses{thm:stinespring_rectangular, thm:kraus_rect_rank_minimal,
        lem:exists_unitary_mul_eq_of_gram_eq}
    Choose the Stinespring dilation space to have dimension $r=dd'$. The
    bound $\rank(\tau_T)\leq dd'$ and the Schr\"odinger form of
    Theorem~\ref{thm:stinespring_rectangular} give an isometry
    $V:\C^d\to\C^{d'}\otimes\C^{dd'}$ satisfying
    $T(\rho)=\tr_{\C^{dd'}}(V\rho V^\dagger)$. Reorder the codomain as
    $\C^d\otimes\C^{d'}\otimes\C^{d'}$, with the retained output factor
    last. Choose a normalized $\varphi\in\C^{d'}\otimes\C^{d'}$ and let
    $W_\varphi x=x\otimes\varphi$. Since the reordered Stinespring map and
    $W_\varphi$ are isometries, their Gram matrices agree. By
    Lemma~\ref{lem:exists_unitary_mul_eq_of_gram_eq}, a unitary $U$ on the
    common $d(d')^2$-dimensional space satisfies $V=UW_\varphi$. Substitution,
    together with
    $W_\varphi\rho W_\varphi^\dagger
      =\rho\otimes\ketbra{\varphi}{\varphi}$,
    gives (\ref{eq:representations_open_system_rectangular_unitary}).
\end{proof}

%% =========================================================
\section{POVMs and Naimark dilation}
\label{sec:povm_naimark}
%% =========================================================

\begin{definition}[Positive operator-valued measure]\label{def:povm}
    \lean{POVM}
    \leanok
    A \emph{positive operator-valued measure} with $n$ outcomes on $\C^D$
    is a family $\{E_i\}_{i=0}^{n-1}$ of positive semidefinite
    operators on $\C^D$ satisfying the resolution of identity
    \begin{align}
        \sum_{i=0}^{n-1}E_i
        &=\Id_D.
        \label{eq:representations_povm_resolution}
    \end{align}
\end{definition}

\begin{lemma}[Unnormalized POVM steering]
    \label{lem:unnormalized_povm_steering}
    \lean{Matrix.exists_povm_of_sum_posSemidef}
    \leanok
    \uses{def:povm, thm:quantum_steering}
    Let $\rho\geq0$ act on a finite-dimensional Hilbert space, and let
    $C:\C^r\to\mathcal H$ satisfy $CC^\dagger=\rho$. If
    $\rho=\sum_{i=0}^{n-1}\sigma_i$, where $n\geq1$ and every
    $\sigma_i\geq0$, then there is a POVM $\{P_i\}_{i=0}^{n-1}$ on
    $\C^r$ such that
    \begin{align}
        \sigma_i
        &=C P_i^{\mathsf T}C^\dagger
        \label{eq:representations_unnormalized_steering}
    \end{align}
    for every $i$.
\end{lemma}

\begin{proof}\leanok
    \uses{thm:quantum_steering}
    Apply the generalized-inverse construction from quantum steering to the
    positive summands $\sigma_i$. Their supports lie in the support of
    $\rho$. The resulting effects sum to the projection onto the part of
    $\C^r$ seen by $C$; add the complementary projection to one effect.
    This preserves (\ref{eq:representations_unnormalized_steering}) and makes
    the effects sum to the identity.
\end{proof}

\begin{theorem}[Environment-induced instruments, supplied Stinespring form]
    \label{thm:environment_induced_instruments_stinespring}
    \lean{Channel.exists_environment_povm_of_sum_eq_stinespring}
    \leanok
    \uses{def:povm, def:kraus_rank_rect, thm:choi_rect_cp,
        lem:unnormalized_povm_steering}
    Let $d\geq1$, let $T:\MN{d}\to\MN{d'}$ be a quantum channel, and let
    $V:\C^d\to\C^{d'}\otimes\C^r$ be a supplied Stinespring matrix such
    that
    \begin{align}
        T(\rho)&=\tr_{\C^r}(V\rho V^\dagger).
        \label{eq:representations_environment_instruments_stinespring}
    \end{align}
    Suppose $n\geq1$ and
    $T=\sum_{i=0}^{n-1}T_i$, where every
    $T_i:\MN{d}\to\MN{d'}$ is completely positive. Then there is a POVM
    $\{P_i\}_{i=0}^{n-1}$ on $\C^r$ such that, for every $i$ and $\rho$,
    \begin{align}
        T_i(\rho)
        &=\tr_{\C^r}\!\left[(\Id_{d'}\otimes P_i)
          V\rho V^\dagger\right].
        \label{eq:representations_environment_instruments_reduced}
    \end{align}
    If $k_i$ is the Kraus rank, equivalently the Choi rank, of $T_i$, then
    $k_i\leq\rank(P_i)$.
\end{theorem}

\begin{proof}\leanok
    \uses{lem:unnormalized_povm_steering, thm:choi_rect_cp,
        def:kraus_rank_rect}
    The supplied Stinespring matrix gives a purification $C$ of the Choi
    matrix $\tau_T$. Complete positivity and
    $T=\sum_iT_i$ give a positive decomposition
    $\tau_T=\sum_i\tau_{T_i}$. Lemma~\ref{lem:unnormalized_povm_steering}
    yields a POVM satisfying
    $\tau_{T_i}=CP_i^{\mathsf T}C^\dagger$. The Choi identity for inserting
    an environment effect gives
    (\ref{eq:representations_environment_instruments_reduced}). Finally,
    \begin{align}
        k_i=\rank(\tau_{T_i})
        =\rank(CP_i^{\mathsf T}C^\dagger)
        \leq\rank(P_i).
        \notag
    \end{align}
\end{proof}

\begin{proposition}[Environment-induced instruments]
    \label{prop:environment_induced_instruments}
    \lean{Channel.exists_environment_povm_of_sum_eq_openSystem}
    \leanok
    \uses{def:povm, def:kraus_rank_rect,
        thm:environment_induced_instruments_stinespring}
    Let $d\geq1$ and let $T:\MN{d}\to\MN{d'}$ be a quantum channel.
    Suppose a normalized vector $\varphi\in\C^{d'}\otimes\C^{d'}$ and a
    unitary identification
    \begin{align}
        U:\C^d\otimes(\C^{d'}\otimes\C^{d'})
        &\longrightarrow \C^{dd'}\otimes\C^{d'}
        \notag
    \end{align}
    are supplied, so both sides have dimension $d(d')^2$, and suppose that
    the system-plus-environment representation from Equation~(2.14) holds:
    \begin{align}
        T(\rho)
        &=\tr_{\C^{dd'}}\!\left[
          U\bigl(\rho\otimes\ketbra{\varphi}{\varphi}\bigr)U^\dagger
          \right].
        \label{eq:representations_environment_instruments_open_system}
    \end{align}
    If $n\geq1$ and $T=\sum_{i=0}^{n-1}T_i$ is a decomposition into
    completely positive maps $T_i:\MN{d}\to\MN{d'}$, then there is a POVM
    $\{P_i\}_{i=0}^{n-1}\subset\MN{dd'}$ such that
    \begin{align}
        T_i(\rho)
        &=\tr_{\C^{dd'}}\!\left[
          (P_i\otimes\Id_{d'})
          U\bigl(\rho\otimes\ketbra{\varphi}{\varphi}\bigr)U^\dagger
          \right]
        \label{eq:representations_environment_induced_instruments}
    \end{align}
    for every $i$ and $\rho$. Moreover, the Kraus rank $k_i$ of $T_i$
    satisfies $k_i\leq\rank(P_i)$. This is
    \cite[Proposition~(Environment induced instruments), Equation~(2.15)]
    {Wolf2012Quantum}.
\end{proposition}

\begin{proof}\leanok
    \uses{thm:environment_induced_instruments_stinespring}
    Restrict $U$ to inputs of the form $x\otimes\varphi$ and reorder the
    output factors. This gives a supplied Stinespring matrix
    $V:\C^d\to\C^{d'}\otimes\C^{dd'}$ satisfying
    (\ref{eq:representations_environment_instruments_stinespring}). Apply
    Theorem~\ref{thm:environment_induced_instruments_stinespring} and rewrite
    its effect-insertion formula in the original environment-first tensor
    order. The result is
    (\ref{eq:representations_environment_induced_instruments}), with the same
    rank bound.
\end{proof}

\begin{definition}[Naimark Kraus square roots]\label{def:naimark_kraus}
    \lean{POVM.naimarkKraus, POVM.naimarkKraus_spec}
    \leanok
    \uses{def:povm}
    For a POVM $\{E_i\}$, each effect $E_i\ge0$ admits a square-root
    factorisation $E_i=M_i^\dagger M_i$ with $M_i\in\MN{D}$. The
    operators $M_i$ are the \emph{Naimark Kraus square roots}.
\end{definition}

\begin{theorem}[Naimark square roots satisfy the Kraus normalization]
    \label{lem:naimark_square_root_normalization}
    \lean{POVM.sum_naimarkKraus_conjTranspose_mul}
    \leanok
    \uses{def:naimark_kraus, def:povm}
    For a POVM $\{E_i\}$ with square roots $E_i=M_i^\dagger M_i$,
    one has $\sum_iM_i^\dagger M_i=\Id$.
\end{theorem}

\begin{proof}\leanok
    Sum the identities $E_i=M_i^\dagger M_i$ and use the defining
    resolution of the identity for the POVM.
\end{proof}

\begin{definition}[Naimark isometry]\label{def:naimark_isometry}
    \lean{POVM.naimarkIsometry}
    \leanok
    \uses{def:naimark_kraus, def:stinespring_v}
    The \emph{Naimark isometry} $V:\C^D\to\C^D\otimes\C^n$ of a POVM
    is the Stinespring-type construction
    \begin{align}
        V
        &=\sum_{i=0}^{n-1}M_i\otimes\ket{i},
        \notag\\
        V_{(a,i),k}
        &=(M_i)_{a,k}.
        \label{eq:representations_naimark_isometry}
    \end{align}
\end{definition}

\begin{definition}[Naimark projective measurement]\label{def:naimark_projection}
    \lean{POVM.naimarkProjection, POVM.naimarkProjection_apply}
    \leanok
    The \emph{Naimark projectors} on $\C^D\otimes\C^n$ are
    \begin{align}
        P_i
        &=\Id_D\otimes\ketbra{i}{i},
        \notag\\
        (P_i)_{(a,b),(c,d)}
        &=\delta_{a,c}\delta_{b,i}\delta_{d,i}.
        \label{eq:representations_naimark_projectors}
    \end{align}
\end{definition}

\begin{theorem}[Naimark isometry condition]\label{thm:naimark_isometry_isometry}
    \lean{POVM.naimarkIsometry_isometry}
    \leanok
    \uses{def:naimark_isometry, def:povm}
    The Naimark isometry satisfies $V^\dagger V=\Id_D$.
\end{theorem}

\begin{proof}
    \leanok
    The Stinespring Gram identity gives
    $V^\dagger V=\sum_iM_i^\dagger M_i=\sum_iE_i=\Id_D$
    by the resolution of identity.
\end{proof}

\begin{theorem}[Naimark projectors form a projective measurement]
    \label{thm:naimark_projection_axioms}
    \lean{POVM.naimarkProjection_mul_self,
        POVM.naimarkProjection_hermitian,
        POVM.naimarkProjection_orthogonal,
        POVM.naimarkProjection_sum_eq_one}
    \leanok
    \uses{def:naimark_projection}
    The Naimark projectors satisfy
    \begin{align}
        P_i^2
        &=P_i,
        \notag\\
        P_i^\dagger
        &=P_i,
        \notag\\
        i\ne j
        &\Rightarrow P_iP_j=0,
        \notag\\
        \sum_iP_i
        &=\Id_{\C^D\otimes\C^n}.
        \notag
    \end{align}
\end{theorem}

\begin{proof}
    \leanok
    Each identity follows from direct entrywise computation using the
    delta-function form in (\ref{eq:representations_naimark_projectors}).
\end{proof}

\begin{theorem}[Naimark dilation]\label{thm:naimark_recovers_povm}
    \lean{POVM.naimark_recovers_povm}
    \leanok
    \uses{def:povm, def:naimark_isometry, def:naimark_projection,
        thm:naimark_isometry_isometry, thm:naimark_projection_axioms}
    Every POVM $\{E_i\}_{i=0}^{n-1}$ arises as a projective measurement on a
    dilation: for the isometry $V$ and projectors $P_i$ above,
    $E_i=V^\dagger P_iV$.
\end{theorem}

\begin{proof}
    \leanok
    Computing entrywise,
    \begin{align}
        (V^\dagger P_iV)_{k,\ell}
        &=\sum_a\overline{(M_i)_{a,k}}(M_i)_{a,\ell}
          =(M_i^\dagger M_i)_{k,\ell}
          =(E_i)_{k,\ell}.
        \notag
    \end{align}
\end{proof}

\begin{theorem}[Existential Naimark dilation]\label{thm:exists_naimark_dilation}
    \lean{POVM.exists_naimark_dilation}
    \leanok
    \uses{thm:naimark_recovers_povm}
    For every POVM $\{E_i\}$ on $\C^D$ there exist a dilation dimension $r$,
    an isometry $V:\C^D\to\C^D\otimes\C^r$, and a projective
    measurement $\{P_i\}$ on the dilation satisfying $E_i=V^\dagger P_iV$.
\end{theorem}

\begin{proof}
    \leanok
    Take $r=n$ and the explicit witnesses $V=\sum_iM_i\otimes\ket{i}$
    and $P_i=\Id_D\otimes\ketbra{i}{i}$.
\end{proof}

\begin{definition}[Naimark dilation as a structure]
    \label{def:is_naimark_dilation}
    \lean{POVM.IsNaimarkDilation}
    \leanok
    \uses{thm:naimark_isometry_isometry, thm:naimark_projection_axioms,
        thm:naimark_recovers_povm}
    A \emph{Naimark dilation} of a POVM $\{E_i\}$ consists of an isometry
    $V$ into a larger Hilbert space together with a projective measurement
    $\{P_i\}$ there such that $E_i=V^\dagger P_iV$ for all $i$.
\end{definition}

\begin{theorem}[Canonical dilation satisfies the Naimark dilation axioms]
    \label{thm:is_naimark_dilation_naimark}
    \lean{POVM.isNaimarkDilation_naimark}
    \leanok
    \uses{def:is_naimark_dilation, thm:naimark_recovers_povm}
    The explicit isometry $V=\sum_iM_i\otimes\ket{i}$ and projectors
    $P_i=\Id_D\otimes\ketbra{i}{i}$ satisfy the defining axioms
    of a Naimark dilation.
\end{theorem}

\begin{proof}
    \leanok
    Combine the previously proved identities $V^\dagger V=\Id_D$,
    $P_i^2=P_i$, $P_i^\dagger=P_i$, $P_iP_j=0$ for $i\ne j$,
    $\sum_iP_i=\Id$, and $V^\dagger P_iV=E_i$.
\end{proof}

\begin{theorem}[Concrete uniqueness for canonical projectors]
    \label{thm:naimark_concrete_uniqueness}
    \lean{POVM.exists_isometry_mul_naimarkIsometry_of_recovery}
    \leanok
    \uses{def:naimark_isometry, def:naimark_projection, thm:naimark_recovers_povm}
    Let $V:\C^D\to\C^D\otimes\C^n$ satisfy
    $V^\dagger P_iV=E_i$ for the canonical projectors
    $P_i=\Id_D\otimes\ketbra{i}{i}$.
    Then there exists an isometry $W$ on the dilated space such that
    $V=WV_0$, where $V_0$ is the canonical Naimark isometry of $\{E_i\}$.
\end{theorem}

\begin{proof}
    \leanok
    For each outcome $i$, the $i$-th block of $V$ has Gram matrix $E_i$.
    Comparing with the canonical square root $M_i$ gives
    $V_i^\dagger V_i=M_i^\dagger M_i$, hence $V_i=U_iM_i$ for a unitary
    $U_i$ on $\C^D$. Assemble the $U_i$ block-diagonally to obtain an
    isometry $W=\bigoplus_iU_i$ on $\C^D\otimes\C^n$, and then
    $V=W\sum_iM_i\otimes\ket{i}=WV_0$.
\end{proof}

\begin{theorem}[Sharp rank-one Naimark extension]
    \label{thm:sharp_rank_one_naimark}
    \lean{POVM.exists_orthonormal_basis_restriction,
        POVM.exists_orthonormal_basis_restriction_of_rank_one}
    \leanok
    \uses{def:povm}
    Let $\{\ketbra{\psi_i}{\psi_i}\}_{i=0}^{n-1}$ be a rank-one POVM on
    $\C^d$ with $n$ outcomes, i.e.\
    $\sum_i\ketbra{\psi_i}{\psi_i}=\Id_d$.  Then necessarily $d\le n$,
    and there exists an orthonormal basis
    $\{\phi_i\}_{i=0}^{n-1}$ of $\C^n$ such that each $\psi_i$ is the
    restriction of $\phi_i$ to the first $d$ coordinates.
    Concretely, the rows of a unitary $U\in U(n)$ give the $\phi_i$, and
    $\psi_i$ is recovered as $\psi_i(j)=U_{i,\,j}$ for
    $j=0,\dots,d-1$.

    Unlike \ref{thm:exists_naimark_dilation}, whose ambient Hilbert space
    is $\C^D\otimes\C^n$ of dimension $D\cdot n$, this theorem yields the
    sharp $n$-dimensional ambient space asserted by
    \cite[Theorem ``Neumark's theorem'']{Wolf2012Quantum}.

    The corresponding corollary starts from a positive operator-valued
    measure with an explicit rank-one decomposition.
\end{theorem}

\begin{proof}
    \leanok
    \uses{def:povm, lem:exists_unitary_mul_eq_of_gram_eq}
    Assemble the $n\times d$ matrix $\Psi_{i,j}:=\psi_i(j)$.  From the
    resolution of identity one obtains $\Psi^\dagger\Psi=\Id_d$.
    Thus $\Psi$ is an isometry, which implies $d\le n$.
    Let $J$ be the $n\times d$ inclusion matrix
    $J_{i,j}=1$ if $i$ equals the embedding of $j$ into $\Fin{n}$ and $0$
    otherwise; for $d\le n$, $J^\dagger J=\Id_d$.
    By Lemma~\ref{lem:exists_unitary_mul_eq_of_gram_eq}, there exists a
    unitary $U\in U(n)$ with $\Psi=UJ$, hence
    $\psi_i(j)=U_{i,\iota(j)}$ where $\iota:\Fin{d}\to\Fin{n}$ is the
    natural inclusion.  The rows of $U$ are the required orthonormal
    basis.  The corollary extracts the vectors from a POVM whose
    effects are given in rank-one form and applies the main theorem.
\end{proof}

\begin{definition}[POVM from a PSD resolution of identity on a dilation]
    \label{def:povm_of_psd_resolution_of_identity}
    \lean{POVM.ofPSDResolutionOfIdentity}
    \leanok
    Given an isometry $V:\C^D\to\C^{d'}$ with $V^\dagger V=\Id_D$
    and a family $\{P_i\}_{i=0}^{n-1}$ of positive semidefinite operators on
    $\C^{d'}$ summing to the identity, the pulled-back operators
    $E_i:=V^\dagger P_iV$ form a POVM. In particular, any projective
    measurement on the dilation (a special case of a PSD resolution of
    identity) pulls back to a POVM.
\end{definition}

\begin{definition}[Quantum instrument]\label{def:instrument}
    \lean{Instrument}
    \leanok
    \uses{def:cp_map, def:trace_preserving}
    A \emph{quantum instrument} with $n$ outcomes is a family
    $\{\Phi_i\}_{i=0}^{n-1}$ of completely positive maps on $\MN{D}$ whose
    sum $\sum_i\Phi_i$ is trace-preserving.
\end{definition}

\begin{definition}[Instrument operations]\label{def:instrument_operations}
    \lean{Instrument.total, Instrument.update, Instrument.probability,
        Instrument.posteriorState}
    \leanok
    \uses{def:instrument}
    Associated to an instrument are the total channel $\sum_i\Phi_i$, the
    unnormalized update $\rho\mapsto\Phi_i(\rho)$ for each outcome $i$, the
    outcome probability $p_i(\rho)=\tr(\Phi_i(\rho))$, and the normalized
    posterior state $\Phi_i(\rho)/p_i(\rho)$ whenever $p_i(\rho)\neq0$.
\end{definition}

\begin{theorem}[Instrument total map is a channel]
    \label{thm:instrument_total_is_channel}
    \lean{Instrument.total_isChannel}
    \leanok
    \uses{def:instrument, def:channel}
    The total map $\sum_i\Phi_i$ of an instrument is a quantum channel.
\end{theorem}

\begin{proof}
    \leanok
    Complete positivity of the sum follows from closure of CP maps under
    finite addition; trace preservation is built into the definition.
\end{proof}

\begin{theorem}[Conservation of probability for instruments]
    \label{thm:instrument_sum_probability}
    \lean{Instrument.sum_probability}
    \leanok
    \uses{def:instrument}
    For every state $\rho\in\MN{D}$, the outcome probabilities
    $p_i(\rho):=\tr(\Phi_i(\rho))$ sum to $\tr(\rho)$.
\end{theorem}

\begin{proof}
    \leanok
    Linearity of trace and trace preservation of $\sum_i\Phi_i$ give
    \begin{align}
        \sum_i\tr(\Phi_i(\rho))
        &=\tr\left(\left(\sum_i\Phi_i\right)(\rho)\right)
         =\tr(\rho).
        \notag
    \end{align}
\end{proof}

\begin{theorem}[Non-negativity of instrument probabilities]
    \label{thm:instrument_probability_nonneg}
    \lean{Instrument.probability_nonneg}
    \leanok
    \uses{def:instrument}
    If $\rho\in\MN{D}$ is positive semidefinite, then each outcome
    probability $p_i(\rho)=\tr(\Phi_i(\rho))$ is a
    non-negative real number.
\end{theorem}

\begin{proof}
    \leanok
    Each $\Phi_i$ is completely positive, hence positive, so
    $\Phi_i(\rho)$ is positive semidefinite and its trace is a
    non-negative real.
\end{proof}

%% =========================================================
